\chapter{Multivariable Differential Calculus}
\index{Multivariable Calculus}
\index{Partial Derivatives}
\index{Total Differential}
\index{Chain Rule}
\index{Directional Derivative}
\index{Tangent Plane}
\index{Lagrange Multipliers}
\index{Thermodynamic Potentials}

\begin{chapterobjectives}
Upon mastering this chapter, students will be able to:
\begin{enumerate}
    \item Compute partial derivatives and verify Clairaut's theorem for mixed second-order derivatives.
    \item Construct total differentials and distinguish exact differentials from inexact differentials in thermodynamics.
    \item Apply the multivariable chain rule to moving frames and coordinate transformations.
    \item Calculate directional derivatives, normal vectors, and tangent planes to level surfaces.
    \item Locate and classify unconstrained critical points using the Hessian determinant test.
    \item Solve constrained physical optimization problems using the method of Lagrange multipliers.
\end{enumerate}
\end{chapterobjectives}

\section{Partial Derivatives and Clairaut's Symmetry Theorem}
Physical fields in classical mechanics, thermodynamics, and electromagnetism depend on multiple spatial and temporal coordinates $(x, y, z, t)$. For a scalar function $f(x, y)$, the \textbf{partial derivative} with respect to $x$ measures the rate of change while holding $y$ strictly constant:
\begin{equation}
\frac{\partial f}{\partial x} \equiv \lim_{\Delta x \to 0} \frac{f(x + \Delta x, y) - f(x, y)}{\Delta x}
\end{equation}

\begin{maththeorem}[Clairaut's Theorem on Mixed Partials]
If the partial derivatives $\frac{\partial^2 f}{\partial x \partial y}$ and $\frac{\partial^2 f}{\partial y \partial x}$ exist and are continuous on an open disk containing point $(x_0, y_0)$, then the order of differentiation is commutative:
\begin{equation}
\frac{\partial^2 f}{\partial x \partial y} = \frac{\partial^2 f}{\partial y \partial x}
\end{equation}
In thermodynamics, this mathematical symmetry generates the powerful \textbf{Maxwell relations} connecting pressure, volume, temperature, and entropy.
\end{maththeorem}

\section{Total Differentials and Thermodynamic State Functions}
The \textbf{total differential} of a function $f(x, y, z)$ represents the first-order change induced by arbitrary infinitesimal increments $(dx, dy, dz)$:
\begin{equation}
df = \frac{\partial f}{\partial x}dx + \frac{\partial f}{\partial y}dy + \frac{\partial f}{\partial z}dz = \nabla f \cdot d\vec{r}
\end{equation}

\begin{mathdefinition}[Exact vs. Inexact Differentials]
A differential form $\delta Q = P(x, y)dx + Q(x, y)dy$ is \textbf{exact} if it can be written as the total differential of a state function $f(x, y)$, which requires:
\begin{equation}
\frac{\partial P}{\partial y} = \frac{\partial Q}{\partial x}
\end{equation}
In thermodynamics, internal energy $dU = TdS - PdV$ is an exact differential (a state function independent of path), whereas heat $\delta Q$ and work $\delta W$ are inexact differentials whose integrals depend explicitly on the thermodynamic process path.
\end{mathdefinition}

\begin{pitfallbox}[Thermodynamic Notation and Constrained Variables]
In thermodynamics, writing $\frac{\partial U}{\partial T}$ is ambiguous because internal energy can be expressed in terms of $(T, V)$ or $(T, P)$. One must always specify the held-constant variable explicitly using subscripts:
\begin{equation}
\left(\frac{\partial U}{\partial T}\right)_V \neq \left(\frac{\partial U}{\partial T}\right)_P
\end{equation}
\end{pitfallbox}

\section{Directional Derivatives and Tangent Planes}
The directional derivative of $f(x, y, z)$ along a unit direction vector $\hat{u} = u_x\hat{i} + u_y\hat{j} + u_z\hat{k}$ is:
\begin{equation}
D_{\hat{u}} f = \nabla f \cdot \hat{u} = |\nabla f|\cos\theta
\end{equation}
The rate of change is maximized when $\hat{u}$ points along the gradient ($\theta = 0$), where $D_{\hat{u}} f = |\nabla f|$. It vanishes when $\hat{u}$ is tangent to the level surface ($\theta = 90^\circ$).

For a level surface defined by $F(x, y, z) = 0$, the gradient vector $\vec{n} = \nabla F(x_0, y_0, z_0)$ is strictly normal to the tangent plane at $(x_0, y_0, z_0)$. The equation of the \textbf{tangent plane} is:
\begin{equation}
\nabla F(x_0, y_0, z_0) \cdot (\vec{r} - \vec{r}_0) = 0 \iff F_x(x - x_0) + F_y(y - y_0) + F_z(z - z_0) = 0
\end{equation}

\section{Constrained Optimization: The Method of Lagrange Multipliers}
Physical systems frequently settle into equilibrium states that minimize potential energy or maximize entropy subject to geometric or conservation constraints.

\begin{maththeorem}[Lagrange Multiplier Principle]
To find the extrema of an objective function $f(x, y, z)$ subject to an equality constraint $g(x, y, z) = 0$, the gradients of $f$ and $g$ must be parallel at the extremum point:
\begin{equation}
\nabla f = \lambda \nabla g
\end{equation}
where $\lambda$ is the \textbf{Lagrange multiplier}. Along with the constraint $g(x, y, z) = 0$, this yields four equations in four unknowns $(x, y, z, \lambda)$.
\end{maththeorem}

\section{Worked Examples and Physical Applications}

\begin{psctexample}[Minimizing Electrostatic Potential Energy on a Constraint Surface]
A charged particle moves on the elliptical cylinder constraint surface $g(x, y) = x^2 + 4y^2 - 4 = 0$ in the presence of an electrostatic potential $V(x, y) = x^2 + y^2 - 2x$. Determine the location of the constrained equilibrium point that minimizes potential energy.

\textbf{\color{rbruNavy}Picture / Conceptualization:}
The particle is physically confined to an elliptical wire loop. The stable equilibrium position corresponds to the lowest point of the potential energy landscape on this ellipse.

\textbf{\color{rbruNavy}Strategy / Governing Laws:}
Apply the Lagrange multiplier condition $\nabla V = \lambda \nabla g$ with the ellipse constraint $g(x, y) = 0$.

\textbf{\color{rbruNavy}Calculation / Analytical Solution:}
1. Compute gradients:
\begin{align}
\nabla V &= (2x - 2)\hat{i} + 2y\hat{j} \\
\nabla g &= 2x\hat{i} + 8y\hat{j}
\end{align}
2. Set up component equations:
\begin{align}
2x - 2 &= 2\lambda x \implies x(1 - \lambda) = 1 \implies x = \frac{1}{1 - \lambda} \\
2y &= 8\lambda y \implies y(1 - 4\lambda) = 0
\end{align}
3. Case Analysis:
\begin{itemize}
    \item \textbf{Case A:} $y = 0$. Substituting into constraint $x^2 + 4(0) = 4 \implies x = \pm 2$.
    \begin{itemize}
        \item At $(2, 0)$: $V(2, 0) = 4 + 0 - 4 = 0$.
        \item At $(-2, 0)$: $V(-2, 0) = 4 + 0 + 4 = 8$.
    \end{itemize}
    \item \textbf{Case B:} $1 - 4\lambda = 0 \implies \lambda = 1/4$.
    Then $x = \frac{1}{1 - 1/4} = \frac{4}{3}$.
    From constraint: $(4/3)^2 + 4y^2 = 4 \implies \frac{16}{9} + 4y^2 = 4 \implies 4y^2 = \frac{20}{9} \implies y = \pm\frac{\sqrt{5}}{3}$.
    Evaluating potential:
    \begin{equation}
    V\left(\frac{4}{3}, \pm\frac{\sqrt{5}}{3}\right) = \left(\frac{4}{3}\right)^2 + \frac{5}{9} - 2\left(\frac{4}{3}\right) = \frac{16 + 5 - 24}{9} = -\frac{3}{9} = -\frac{1}{3}
    \end{equation}
\end{itemize}

\textbf{\color{rbruNavy}Think / Physical Insights:}
Comparing the values: $V(4/3, \pm\sqrt{5}/3) = -1/3$ is strictly lower than $V(2, 0) = 0$ and $V(-2, 0) = 8$. Thus, there are two symmetric stable equilibrium points at $(4/3, \pm\sqrt{5}/3)$.
\qedexam
\end{psctexample}

\begin{pythonbox}[Lagrange Multiplier Contour Plot in Python]
import numpy as np

# Verify potential values on the ellipse x = 2*cos(t), y = sin(t)
t = np.linspace(0, 2*np.pi, 1000)
x = 2 * np.cos(t)
y = np.sin(t)
V = x**2 + y**2 - 2*x

min_idx = np.argmin(V)
print(f"Numerical Minimum V = {V[min_idx]:.4f} at x = {x[min_idx]:.4f}, y = {y[min_idx]:.4f}")
print(f"Analytical Minimum: -1/3 = {-1/3:.4f}, x = 4/3 = {4/3:.4f}, y = sqrt(5)/3 = {np.sqrt(5)/3:.4f}")
\end{pythonbox}

\section{Summary of Key Formulas}
\begin{table}[H]
\centering
\small
\begin{tabularx}{\textwidth}{llX}
\toprule
\textbf{Theorem / Concept} & \textbf{Formula} & \textbf{Physical Application} \\
\midrule
Clairaut's Symmetry & $\frac{\partial^2 f}{\partial x \partial y} = \frac{\partial^2 f}{\partial y \partial x}$ & Maxwell relations in thermodynamics, integrability of potentials \\
Total Differential & $df = \sum_i \frac{\partial f}{\partial x_i}dx_i$ & First law of thermodynamics, propagation of experimental errors \\
Directional Derivative & $D_{\hat{u}} f = \nabla f \cdot \hat{u}$ & Temperature gradients in heat conduction, pressure gradients in meteorology \\
Lagrange Multipliers & $\nabla f = \sum_k \lambda_k \nabla g_k$ & Microcanonical entropy maximization, constraint forces in D'Alembert mechanics \\
\bottomrule
\end{tabularx}
\caption{Key relations in multivariable differential calculus and thermodynamics.}
\label{tab:ch06_summary}
\end{table}

\section{Chapter Exercises}
\begin{enumerate}
    \item \textbf{Thermodynamic Maxwell Relation:} From the Helmholtz free energy $F = U - TS$ with $dF = -SdT - PdV$, prove that $\left(\frac{\partial S}{\partial V}\right)_T = \left(\frac{\partial P}{\partial T}\right)_V$.
    \item \textbf{Tangent Plane to an Ellipsoid:} Find the equation of the tangent plane to the ellipsoid $\frac{x^2}{a^2} + \frac{y^2}{b^2} + \frac{z^2}{c^2} = 1$ at an arbitrary point $(x_0, y_0, z_0)$.
    \item \textbf{Maximum Area with Fixed Perimeter:} Use Lagrange multipliers to prove that among all rectangles of fixed perimeter $P$, the square encloses the maximum area.
\end{enumerate}
